% Fixture: math-heavy
% Dense inline mathematics, numbered displays (equation, align, cases), a
% matrix, and definition/theorem/lemma/proof environments, on A4. Each page
% ends with an explicit \clearpage so page contents are fixed by this source.
\documentclass[10pt]{article}
% Shared settings for every Tourmaline fixture. \input this right after
% \documentclass.
%
% Reproducible output: no creation date, no random trailer id, no pdfTeX
% banner, so rebuilding with the same TeX Live gives byte-identical PDFs.
\pdfinfoomitdate=1
\pdftrailerid{}
\pdfsuppressptexinfo=-1
% Map every glyph (including the fi/fl/ffi ligatures) to Unicode so that
% text extraction does not depend on font-specific guesswork.
\input{glyphtounicode}
\pdfgentounicode=1

\usepackage[a4paper,margin=25mm]{geometry}
\usepackage{amsmath,amssymb,amsthm}

\newtheorem{theorem}{Theorem}
\newtheorem{lemma}[theorem]{Lemma}
\theoremstyle{definition}
\newtheorem{definition}[theorem]{Definition}

\title{A Note on Contractive Update Rules}
\author{Emil Stand-In}
\date{}

\begin{document}
\maketitle

% ------------------------------------------------------------------- page 1
\section{Setting}

Throughout this note we work in a finite dimensional real vector space with
the usual inner product. Let $V=\mathbb{R}^n$ with norm $\|x\|=\sqrt{\langle
x,x\rangle}$, and let $f\colon V\to V$ be a map with $f(0)=0$. For
$x_0\in V$ we define the sequence $x_{k+1}=f(x_k)$ for $k\ge 0$, and we ask
when $x_k\to 0$ as $k\to\infty$ for every choice of $x_0$ with
$\|x_0\|\le r$ and some fixed $r>0$.

\begin{definition}
A map $f$ is called contractive on a set $S\subseteq V$ if there is a
constant $0\le q<1$ such that $\|f(x)-f(y)\|\le q\,\|x-y\|$ for all
$x,y\in S$.
\end{definition}

The simplest contractive maps are linear. If $f(x)=Ax$ for a matrix $A$
then
\begin{equation}
  \|x_k\| = \|A^k x_0\| \le \|A\|^k\,\|x_0\| ,
  \label{eq:linear}
\end{equation}
so the iteration converges whenever the operator norm of the matrix is
smaller than one. A concrete example in two dimensions is the matrix
\begin{equation}
  A = \begin{pmatrix} \tfrac12 & \tfrac14 \\[2pt] 0 & \tfrac13 \end{pmatrix},
  \qquad
  A^k = \begin{pmatrix} 2^{-k} & c_k \\ 0 & 3^{-k} \end{pmatrix},
  \label{eq:matrix}
\end{equation}
where $c_k = \tfrac32\,(2^{-k}-3^{-k})$ for all $k\ge 0$.

\begin{theorem}\label{thm:main}
Let $f$ be contractive on the closed ball $B_r=\{x:\|x\|\le r\}$ with
constant $q$, and suppose $f(B_r)\subseteq B_r$. Then for every $x_0\in B_r$
the sequence of iterates satisfies $\|x_k\|\le q^k r$, and in particular it
converges to the origin.
\end{theorem}

\clearpage
% ------------------------------------------------------------------- page 2
\section{Proofs}

The proof of the main theorem rests on a simple estimate that we record
separately, because it is used again in the next section.

\begin{lemma}\label{lem:step}
Under the assumptions of Theorem~\ref{thm:main}, for all $k\ge 0$ we have
\begin{align}
  \|x_{k+1}\| &= \|f(x_k)-f(0)\|   \label{eq:step1}\\
              &\le q\,\|x_k-0\|    \label{eq:step2}\\
              &= q\,\|x_k\| .      \label{eq:step3}
\end{align}
\end{lemma}

\begin{proof}
The first line uses $f(0)=0$, the second is the contraction property on
$B_r$, which applies because both $x_k$ and $0$ lie in the ball, and the
third is immediate.
\end{proof}

\begin{proof}[Proof of Theorem~\ref{thm:main}]
By induction on $k$ and Lemma~\ref{lem:step}, $\|x_k\|\le q^k\|x_0\|\le
q^k r$. Since $0\le q<1$ we have $q^k\to 0$, hence $x_k\to 0$.
\end{proof}

For the example in \eqref{eq:matrix}, the operator norm satisfies
$\|A\|_2\approx 0.557$, which gives the explicit rate
\begin{equation}
  \|x_k\| \le (0.557)^k\,\|x_0\|
  \quad\text{and}\quad
  \sum_{k=0}^{\infty}\|x_k\| \le \frac{\|x_0\|}{1-0.557} .
  \label{eq:rate}
\end{equation}

\section{A piecewise rule}

Consider now the piecewise definition
\begin{equation}
  g(x) =
  \begin{cases}
    \tfrac12 x & \text{if } \|x\|\le 1,\\[2pt]
    \dfrac{x}{2\|x\|^2} & \text{otherwise.}
  \end{cases}
  \label{eq:piecewise}
\end{equation}
Inside the unit ball this map halves every vector, and outside it the image
lies strictly inside the ball, so every orbit enters the ball after one step
and then converges at the rate given by \eqref{eq:linear} with $q=\tfrac12$.

\end{document}
